%%%%%%%%%%%%%%%%%%%%%%%%%%%%% Define Article %%%%%%%%%%%%%%%%%%%%%%%%%%%%%%%%%%
\documentclass{article}
%%%%%%%%%%%%%%%%%%%%%%%%%%%%%%%%%%%%%%%%%%%%%%%%%%%%%%%%%%%%%%%%%%%%%%%%%%%%%%%

%%%%%%%%%%%%%%%%%%%%%%%%%%%%% Using Packages %%%%%%%%%%%%%%%%%%%%%%%%%%%%%%%%%%
\usepackage{geometry}
\usepackage{graphicx}
\usepackage{amssymb}
\usepackage{amsmath}
\usepackage{amsthm}
\usepackage{empheq}
\usepackage{mdframed}
\usepackage{booktabs}
\usepackage{lipsum}
\usepackage{graphicx}
\usepackage{color}
\usepackage{psfrag}
\usepackage{pgfplots}
\usepackage{bm}
%%%%%%%%%%%%%%%%%%%%%%%%%%%%%%%%%%%%%%%%%%%%%%%%%%%%%%%%%%%%%%%%%%%%%%%%%%%%%%%

% Other Settings

%%%%%%%%%%%%%%%%%%%%%%%%%% Page Setting %%%%%%%%%%%%%%%%%%%%%%%%%%%%%%%%%%%%%%%
\geometry{a4paper}

%%%%%%%%%%%%%%%%%%%%%%%%%% Define some useful colors %%%%%%%%%%%%%%%%%%%%%%%%%%
\definecolor{ocre}{RGB}{243,102,25}
\definecolor{mygray}{RGB}{243,243,244}
\definecolor{deepGreen}{RGB}{26,111,0}
\definecolor{shallowGreen}{RGB}{235,255,255}
\definecolor{deepBlue}{RGB}{61,124,222}
\definecolor{shallowBlue}{RGB}{235,249,255}
%%%%%%%%%%%%%%%%%%%%%%%%%%%%%%%%%%%%%%%%%%%%%%%%%%%%%%%%%%%%%%%%%%%%%%%%%%%%%%%

%%%%%%%%%%%%%%%%%%%%%%%%%% Define an orangebox command %%%%%%%%%%%%%%%%%%%%%%%%
\newcommand\orangebox[1]{\fcolorbox{ocre}{mygray}{\hspace{1em}#1\hspace{1em}}}
%%%%%%%%%%%%%%%%%%%%%%%%%%%%%%%%%%%%%%%%%%%%%%%%%%%%%%%%%%%%%%%%%%%%%%%%%%%%%%%

%%%%%%%%%%%%%%%%%%%%%%%%%%%% English Environments %%%%%%%%%%%%%%%%%%%%%%%%%%%%%
\newtheoremstyle{mytheoremstyle}{3pt}{3pt}{\normalfont}{0cm}{\rmfamily\bfseries}{}{1em}{{\color{black}\thmname{#1}~\thmnumber{#2}}\thmnote{\,--\,#3}}
\newtheoremstyle{myproblemstyle}{3pt}{3pt}{\normalfont}{0cm}{\rmfamily\bfseries}{}{1em}{{\color{black}\thmname{#1}~\thmnumber{#2}}\thmnote{\,--\,#3}}
\theoremstyle{mytheoremstyle}
\newmdtheoremenv[linewidth=1pt,backgroundcolor=shallowGreen,linecolor=deepGreen,leftmargin=0pt,innerleftmargin=20pt,innerrightmargin=20pt,]{theorem}{Theorem}[section]
\theoremstyle{mytheoremstyle}
\newmdtheoremenv[linewidth=1pt,backgroundcolor=shallowBlue,linecolor=deepBlue,leftmargin=0pt,innerleftmargin=20pt,innerrightmargin=20pt,]{definition}{Definition}[section]
\theoremstyle{myproblemstyle}
\newmdtheoremenv[linecolor=black,leftmargin=0pt,innerleftmargin=10pt,innerrightmargin=10pt,]{problem}{Problem}[section]
\theoremstyle{myproblemstyle}
\newmdtheoremenv[linecolor=black,leftmargin=0pt,innerleftmargin=10pt,innerrightmargin=10pt,]{example}{Example}[section]
%%%%%%%%%%%%%%%%%%%%%%%%%%%%%%%%%%%%%%%%%%%%%%%%%%%%%%%%%%%%%%%%%%%%%%%%%%%%%%%

%%%%%%%%%%%%%%%%%%%%%%%%%%%%%%% Plotting Settings %%%%%%%%%%%%%%%%%%%%%%%%%%%%%
\usepgfplotslibrary{colorbrewer}
\pgfplotsset{width=8cm,compat=1.9}
\parskip=3pt
\parindent=0pt
%%%%%%%%%%%%%%%%%%%%%%%%%%%%%%%%%%%%%%%%%%%%%%%%%%%%%%%%%%%%%%%%%%%%%%%%%%%%%%%

%%%%%%%%%%%%%%%%%%%%%%%%%%%%%%% Title & Author %%%%%%%%%%%%%%%%%%%%%%%%%%%%%%%%
\title{Diophantine Equations}
\author{Alston Yam}
\date{}
%%%%%%%%%%%%%%%%%%%%%%%%%%%%%%%%%%%%%%%%%%%%%%%%%%%%%%%%%%%%%%%%%%%%%%%%%%%%%%%

\begin{document}
    \maketitle

    \section{Introduction}
    Diophantine equations are equations that ask for $\textbf{integer}$ solutions. This can come up in many places, not just in Olympiad but also normal maths too. Here are a few ways to solve them.

    \section{Factorisation}
    For solutions in integers, if we can factor the two sides so that one side gets to a constant, the rest of the problem will solve itself by simple case work. Here are a few factorising techniques, motivated by examples:

    \begin{example}
        Solve 
        \[x^2 = y^2 + 7\]
    \end{example}

    And so this is very doable by simply factoring into $(x-y)(x+y) = 7$ and by doing case work.

    \subsection{Sophie Germain Identity — who said we can't factor sum of squares?}

    Consider the expression 
    \[x^4 + 4y^4\]

    It's a useful trick to recognise that this expression is in fact factorisable.

    \[x^4 + 4y^4 = x^4 + 4y^4 + 4x^2y^2 - 4x^2y^2 = (x^2 + 2y^2) - 4x^2y^2 = (x^2 - 2xy + 2y^2)(x^2 + 2xy + 2y^2)\]

    \begin{problem}[Example Problem 1]
        Find all solutions in integers for the equation
        \[xy = x + y + 4\]
    \end{problem}

    \begin{problem}[Example Problem 2]
        Find all solutions in integers for the equation
        \[x^2 + y^2 = x + y + 2\]
    \end{problem}

    \section{Using Mod}
    Specifically, we can analyse the outputs of numbers when we raise them to some power. This is sometimes also know as the quadratic residue (for 2nd powers).

    This leads naturally to Fermat's Little Theorem. As a reminder:

    \begin{theorem}[Fermat's Little Theorem]
        \[a^{p-1} \equiv 1 \pmod{p}\]
    \end{theorem}

    However, problems can be solved with a clever mod.

    \begin{example}[RMO 2017/2]
        Show that the equation 
        \[a^3 + (a+1)^3 + \cdots + (a+6)^3 = b^4 + (b+1)^4\]

        Has no solutions in $a$, $b$.
    \end{example}

    The key observation here is to just use mod $7$. The specific details are left as an exercise for the reader.

    \section{Infinite Descent}
    The idea of using the technique of infinite descent is that suppose we find a solution $x_0$. Then by some manipulation, we find another solution $x_1$. This way, we can generate infinite solutions.

    If we have some additional requirements such as $x$ have to be positive, then if we prove $x_i > x_{i+1}$ then we would find a contradiction which is good.

    \begin{example}[APMO 2017/1]
        We call a $5$-tuple of integers \textit{arrangeable} if its elements can be labeled $a, b, c, d, e$ in some order so that $a-b+c-d+e=29$. Determine all $2017$-tuples of integers $n_1, n_2, . . . , n_{2017}$ such that if we place them in a circle in clockwise order, then any $5$-tuple of numbers in consecutive positions on the circle is arrangeable.
    \end{example}

    We can replace each value $n_i$ on the circle with $m_i = n_i - 29$. This is consistent with the problem conditions. An added benefit is that we can now see $m_i \equiv m_{i+5} \pmod{2}$, which we can now use to loop around the circle to see that all $m_i$ must have the same parity.

    Now, we can argue that all $m_i$ must be even, by considering we have $$m_i - m_{i+1} + m_{i+2} - m_{i+3} + m_{i+4} = 0.$$

    This means that we can divide all $m_i$ by $2$ and still have a valid arrangement. But this is a contradiction! As we can do this forever, which will eventually force all of the $m_i$s to be 0. This in turn, when placed back into our original problem, corresponds too all values being 29.

    \pagebreak

    \section{Problems}
    \begin{problem}
        Are there integers $a, b$ such that both $a^5b + 3$ and $ab^5 + 3$ are perfect cubes?
    \end{problem}

    \begin{problem}
        Prove that for natural numbers $m$, $n$, 
        \[3^m + 3^n + 1\]
        cannot be a perfect square.
    \end{problem}





    
\end{document}